% =====================================================================
%  3.6 Modellering — fasit
%  Regresjonskonstanter er avrundet. Små avvik mellom verktøy er greit.
% =====================================================================

Konstantene er avrundet. Digitale verktøy regner litt ulikt, særlig for eksponentielle modeller og potensmodeller. Små avvik i konstantene og i $R^2$ er derfor greit. Det viktige er hvilken modell som har høyest $R^2$, ikke det nøyaktige tallet.

\fk{Del 1 — Regresjon med digitale verktøy}
\fo{1.1} a) Lineær: $y$ øker med 4 hver gang. \ b) Eksponentiell: $y$ ganges med 3 hver gang. \ c) Potens: når $x$ dobles, ganges $y$ med 4. Det er $y = 5x^2$.
\fo{1.2} a) $45$: vannmengden øker med 45 liter per minutt. $120$: det er 120 liter i bassenget ved start.
b) $800$: 800 følgere ved start. $1{,}15$: antallet ganges med $1{,}15$ per uke, altså $15\,\%$ vekst per uke.
c) $V(10) = 570$ liter og $N(2) = 800 \cdot 1{,}15^2 = 1\,058$ følgere.
\fo{1.3} a) Ja, punktene ligger nesten på en rett linje. \ b) $f(x) \approx 52{,}0x + 120$ \ c) 52 elbiler per år
\fo{1.4} a) Differanser: $40, 55, 75, 90, 130, 180$ (vokser, ikke lineær). Forhold: omtrent $1{,}33;\ 1{,}34;\ 1{,}35;\ 1{,}31;\ 1{,}34;\ 1{,}35$, nesten like. Eksponentiell.
b) $f(x) \approx 120 \cdot 1{,}34^x$ \ c) $a \approx 120$: omtrent 120 visninger på dag 0. $b \approx 1{,}34$: antall visninger ganges med $1{,}34$ per dag, altså $34\,\%$ vekst per dag. \ d) Omtrent 1\,200 visninger.
\fo{1.5} a) $1{,}42$ og $\frac{2{,}01}{1{,}42} \approx 1{,}42$. Samme faktor når $L$ dobles, så potensmodell. \ b) $T \approx 2{,}01 \cdot L^{0{,}50}$ \ c) $T(3{,}0) \approx 3{,}5$ s \ d) Omtrent $2{,}2$ m.
\fo{1.6} Steg 1–3 er riktige. Feilen er i steg 4: $1{,}2^x$ med $x = 5$ betyr fem faktorer $1{,}2$, ikke $1{,}2 \cdot 5$. Riktig er $f(5) = 50 \cdot 1{,}2^5 \approx 124{,}4$.
\fo{1.7} a) Alltid. Det er definisjonen på en lineær sammenheng: $y = ax + b$ med $a$ lik tillegget.
b) Noen ganger. Bare når punktene ligger nøyaktig på kurven, for eksempel to punkter og en rett linje. Ellers ligger noen punkter ved siden av.
c) Noen ganger, bare når stigningstallet er 0. Ellers er stigningstallet per måned $\frac{1}{12}$ av det per år, fordi det er 12 ganger så mange $x$-enheter.
\fo{1.8} Eksempel: $y = 20, 30, 45, 68, 101$. Jeg ganget med omtrent $1{,}5$ hver gang. Eksponentiell regresjon gir $f(x) \approx 20{,}0 \cdot 1{,}50^x$ med $R^2 \approx 1{,}00$. Lineær gir $R^2 \approx 0{,}95$. Flere svar er riktige, så lenge faktoren er omtrent lik og tydelig større enn 1, for eksempel 1,5 eller mer (eller mindre enn 0,7).
\fo{Sjekk 1} a) Eksponentiell: $y$ ganges med $\frac{1}{2}$ hver gang. \ b) Da er $x = 0$ første året, og konstanten blir verdien ved start. Årstall gir store, meningsløse konstanter. \ c) $3{,}5$: batteriprosenten øker med 3,5 prosentpoeng per minutt. $12$: batteriet var på $12\,\%$ da ladingen startet.

\fk{Del 2 — Vurdere modeller og lese tekster}
\fo{2.1} a) $P(4) = 52$ kr \ b) $P(35) = -10$. En pris kan ikke være negativ, så modellen gjelder ikke så langt ut. \ c) $0 \le x \le 10$ (dataene). Aldri over $x = 30$, der prisen blir 0.
\fo{2.2} Eksempler på gode svar (flere er riktige):
a) «Ifølge modellen blir folketallet omtrent 9\,000 i 2040, hvis utviklingen fortsetter.» b) «Ifølge modellen dobles prisen omtrent hvert femte år for tidsrommet vi har data om.» c) «$R^2 = 0{,}99$ betyr at kurven passer godt med punktene, men modellen kan likevel være feil for situasjonen.» d) $R^2 = 0{,}12$: modellen passer dårlig med punktene. Se på punktplottet og prøv en annen modelltype, og vurder om modellen passer med situasjonen.
\fo{2.3} a) $f(x) \approx -2{,}19x^2 + 14{,}8x + 12{,}3$ \ b) Toppunkt ved $x \approx 3{,}4$, altså omtrent kl. 13.20, med omtrent 37 besøkende. \ c) $f(8) \approx -9{,}5$. Antall besøkende kan ikke være negativt. \ d) $0 \le x \le 7$ (dataene). Modellen blir negativ etter $x \approx 7{,}5$, så den gjelder ikke senere enn det.
\fo{2.4} a) $x = 0, 4, 8$ med $y = 5\,200, 5\,640, 6\,080$. \ b) $\frac{6\,080 - 5\,200}{5\,200} \approx 16{,}9\,\%$. Påstanden «over 20 prosent» stemmer ikke.
c) Økning $110$ per år: $f(x) = 110x + 5\,200$. For 2040 er $x = 25$, så $f(25) = 7\,950$. \ d) Trengs $\frac{8\,500 - 5\,200}{25} = 132$ per år, altså 22 mer per år enn i modellen. Ordføreren må anta at veksten blir større enn hittil. En eksponentiell modell gjennom 2015 og 2023 gir omtrent 8\,480 i 2040, så ordføreren har trolig regnet med fast prosentvis vekst. Dataene har like store tillegg (440 hvert fjerde år), og det peker mot en lineær modell.
\fo{2.5} Steg 1 og 2 er riktige. Første feil er i steg 3: $\frac{60}{2{,}5} = 24$, ikke 25. Riktig år er $2020 + 24 = 2044$. Steg 5 bør også være «ifølge modellen».
\fo{2.6} a) Noen ganger. En andregradsmodell har alltid minst like høy $R^2$ som en lineær modell for de samme dataene, men kan gi urimelige prognoser. For eksempel kan en synkende bruktpris begynne å stige igjen etter bunnpunktet. Modelltypen må passe med situasjonen.
b) Alltid, for store nok $x$. Hver økning i $x$ trekker fra like mye, og til slutt er $y$ under 0. (Modellen er ikke nødvendigvis gyldig så langt.)
c) Aldri. Når $x$ skiftes, endres bare startverdien $a$ (ganges med $b$ opphøyd i skiftet). Fra ett år til neste ganges $y$ fortsatt med samme $b$.
\fo{2.7} a) Lineær: $y \approx 13{,}6x + 98$ med $R^2 \approx 0{,}996$. Eksponentiell: $y \approx 100 \cdot 1{,}11^x$ med $R^2 \approx 0{,}9998$.
b) Lineær: omtrent 302 tusen. Eksponentiell: omtrent 480 tusen.
c) Gode poeng (2–3 holder): begge passer innenfor dataene, så $R^2$ avgjør ikke. Eksponentiell vekst på $11\,\%$ per år kan ikke fortsette, fordi antall husstander er begrenset og markedet metter seg. Lineær er et forsiktigere anslag, men heller ikke sikker. Nevn forbehold og bruk «ifølge modellen».
\fo{2.8} a) $5 \cdot 1{,}015^{200} \approx 98$ millioner. \ b) Landet har begrenset plass og begrenset med ressurser som mat og vann. Veksten kan ikke bli ved med å være $1{,}5\,\%$ i 200 år. \ c) En kurve som stiger først, deretter flater ut mot en øvre grense. Det finnes modeller med øvre grense. Flere gode svar finnes.
\fo{Sjekk 2} a) $0 \le x \le 12{,}5$ \ b) «Ifølge modellen blir det omtrent 500 brukere i 2030, hvis utviklingen fortsetter.» \ c) Høy $R^2$ sier bare at modellen passer med punktene vi har. Langt utenfor dataene kan situasjonen endre seg.

\fk{Ekstra utfordring}
\fo{U1} a) $\frac{95}{12} \approx 7{,}9$ og $\frac{760}{95} = 8{,}0$. Samme faktor når $L$ dobles, så potensmodell. \ b) $W \approx 0{,}0121 \cdot L^{2{,}99}$ \ c) Vekten ganges med omtrent $2^{3} = 8$. \ d) Omtrent 2\,500–2\,600 g. Usikkert, fordi 60 cm er utenfor dataene (10–40 cm) og fisken kanskje ikke vokser etter samme mønster.
\fo{U2} a) $x = 0, 1, 2$ med $y = 50\,000, 90\,000, 162\,000$. \ b) Forholdene er $1{,}8$ og $1{,}8$. Eksponentiell: $f(x) \approx 50\,000 \cdot 1{,}8^x$. \ c) $x = 8$ gir omtrent 5,5 millioner. \ d) Regnestykket stemmer: ifølge modellen nås 5,5 millioner rundt 1. januar 2031 ($x \approx 8{,}0$). Gode poeng: (1) Påstanden står på en knivsegg. Med vekstfaktor $1{,}79$ i stedet for $1{,}8$ gir modellen bare omtrent 5,3 millioner. (2) Modellen bygger på tre tall over to år. (3) Teksten bytter fra «brukere» til «lastet den ned». (4) «Alle» er umulig, og veksten må flate ut.
\fo{U3} a) \lstinline[style=laereverk,numbers=none,frame=none]|nevner = nevner + (x[i] - x_snitt) ** 2| og \lstinline[style=laereverk,numbers=none,frame=none]|b = y_snitt - a * x_snitt|.
b) Utskrift: omtrent $2{,}657$ og $-15{,}295$. Samme som regresjon (avrundet $a \approx 2{,}66$, $b \approx -15{,}3$). \ c) $a$: salget øker med omtrent 2,66 tusen kr per grad. $b$: negativt salg ved 0 grader. Det er meningsløst, fordi 0 grader er langt utenfor dataene (14–29 grader).
\fo{U4} Eksempel: $y = 2, 6, 17, 50, 150$. Når $x$ dobles, ganges $y$ med omtrent 3. Potensregresjon gir $y \approx 2{,}0 \cdot x^{1{,}55}$ med $R^2 \approx 1{,}00$. Eksponentiell passer klart dårligere ($R^2$ godt under 0,9, tallet varierer med verktøyet). Flere svar er riktige, så lenge $y$ ganges med omtrent samme tall, tydelig forskjellig fra 1, hver gang $x$ dobles.
\fo{U5} a) $a \bar x + b = a \bar x + (\bar y - a \bar x) = \bar y$. Linja gir altså $\bar y$ når $x = \bar x$. \ b) Etter forskyvningen er $x_i' = x_i + k$ og $\bar x' = \bar x + k$. Differansene $x_i' - \bar x'$ er de samme som før, så $a$ er uendret. Da blir $b' = \bar y - a(\bar x + k) = b - ak$. Eksempel 2 i del 1: $211{,}7 - 17{,}0 \cdot 2010 \approx -33\,958$.
\fo{U6} a) $y \approx 3{,}54x^2 - 56{,}1x + 361$ med $R^2 \approx 0{,}9998$. \ b) Bunnpunktet ligger ved $x \approx 7{,}9$ (fordi $a > 0$). Etter det stiger parabelen. For 15 år gir modellen omtrent 315 tusen kr, mer enn for en ett år gammel bil (omtrent 309 tusen kr). Urimelig.
c) Eksponentiell: $y \approx 327 \cdot 0{,}89^x$ med $R^2 \approx 0{,}97$. For 15 år gir den omtrent 50–55 tusen kr (avhenger av verktøyet). Tre poeng: (1) Andregradsmodellen er urimelig fordi den stiger igjen. (2) Eksponentiell er minst urimelig, men fanger ikke at prisene flater ut, og avvikene fra punktene er systematiske. (3) Ingen av modellene er sikre 7 år utenfor dataene. Bruk «ifølge modellen».
